\documentclass[a4paper,11pt]{article}
\usepackage[utf8]{inputenc}
\usepackage[T1]{fontenc}
\usepackage[french]{babel}
\usepackage{lmodern}
\usepackage{amsmath,amssymb,amsthm}
\usepackage{mathrsfs}
\usepackage{enumitem}
\usepackage{multicol}

\setlist[enumerate]{label=\textbf{\textcolor{Couleur8}{\arabic*.}}, left=1em}

% Si vous voulez aussi modifier les listes enumerate imbriquées :
\setlist[enumerate,2]{label=\textcolor{Couleur8}{\alph*)}}
\setlist[enumerate,3]{label=\textcolor{Couleur8}{\roman*)}}
\setlist[enumerate,4]{label=\textcolor{Couleur8}{\Alph*)}}
\usepackage{graphicx}
\usepackage{xcolor}
\usepackage{tcolorbox}
\usepackage{geometry}
\usepackage{fancyhdr}
\usepackage{titlesec}
\usepackage{cancel}
\usepackage{ifthen}
\usepackage{tikz}
\usetikzlibrary{arrows.meta}
\usepackage{tkz-tab}
\usepackage{eurosym}

% OPTION PROF/ÉLÈVE
% Décommentez UNE des deux lignes suivantes :
\newboolean{prof}\setboolean{prof}{true}   % Version professeur (avec corrections des devoirs)
\newboolean{prof}\setboolean{prof}{true}  % Version élève (sans corrections des devoirs)

\definecolor{Couleur1}{HTML}{4A5D7A} %Th
\definecolor{Couleur2}{HTML}{A5B5D6} %Prop
\definecolor{Couleur3}{HTML}{A8C68A} %Méthode
\definecolor{Couleur6}{HTML}{8A9CC1} %Def
\definecolor{Couleur7}{HTML}{E6B459} %Rq Voc
\definecolor{Couleur8}{HTML}{D36740} %Titres
\definecolor{correction}{HTML}{D36740} % Couleur pour les corrections
\definecolor{coopmathscorr}{HTML}{F15929} % Couleur corrections coopmaths

% Configuration des marges
\geometry{
	top=2cm,
	bottom=2cm,
	inner=1.5cm,
	outer=1.5cm,
}

% Police
\renewcommand{\familydefault}{\sfdefault}

% Compteurs
\newcounter{seancenum}
\newcounter{seancenumD}
\setcounter{seancenum}{1}
\setcounter{seancenumD}{1}

% Configuration des en-têtes et pieds de page
\pagestyle{fancy}
\fancyhf{}
\ifthenelse{\boolean{prof}}{
    \fancyhead[L]{\textcolor{Couleur8}{\textbf{Corrections Automatismes - Classe de seconde - Version Professeur}}}
}{
    \fancyhead[L]{\textcolor{Couleur8}{\textbf{Corrections Devoirs - Séance 16 - Classe de seconde}}}
}
\fancyhead[R]{\textcolor{Couleur8}{\textbf{2025-2026}}}
\fancyfoot[L]{\textcolor{Couleur8}{Mme Bertail et Mme Pinto}}
\fancyfoot[R]{\textcolor{Couleur8}{\thepage}}
\renewcommand{\headrulewidth}{1pt}
\renewcommand{\headrule}{\hbox to\headwidth{\color{Couleur8}\leaders\hrule height \headrulewidth\hfill}}

% Environnements personnalisés
\newtcolorbox{seance}[1]{
    colback=Couleur6!10,
    colframe=Couleur1!65,
    fonttitle=\bfseries\large,
    title=Correction Séance \theseancenum #1,
    left=5mm,
    right=5mm,
    top=3mm,
    bottom=3mm,
    arc=0mm,
    after=\stepcounter{seancenum}
}

\newtcolorbox{seanceD}[1]{
    colback=Couleur6!5,
    colframe=Couleur6!45,
    fonttitle=\bfseries\large,
    title=Correction Devoirs - Séance \theseancenumD #1,
    left=5mm,
    right=5mm,
    top=3mm,
    bottom=3mm,
    arc=0mm,
    after=\stepcounter{seancenumD}
}

\begin{document}

\setcounter{seancenumD}{16}
\newpage
\begin{seanceD}{} %16

\begin{enumerate}
\item $(-4x-12)(-9x+4)>0$\\


\begin{tikzpicture}[baseline, scale=0.5]
\tkzTabInit[lgt=10,deltacl=0.8,espcl=2.5]{ $x$ / 2.5, $-4x-12$ / 2, $-9x+4$ / 2, $(-4x-12)(-9x+4)$ / 2}{ $-\infty$, $-3$, $\dfrac{4}{9}$, $+\infty$}
\tkzTabLine{  , +, z, -, t, -}
\tkzTabLine{  , +, t, +, z, -}
\tkzTabLine{  , +, z, -, z, +}
\end{tikzpicture}

L'ensemble de solutions de l'inéquation est $S ={\color{correction}\boldsymbol{ \bigg] -\infty ; -3 \bigg[ \cup \bigg] \dfrac{4}{9}; +\infty \bigg[ }}$.

\item $64x^2-80x+25=(8x)^2-2 \times 8x \times 5 + 5^2={\color{correction}\boldsymbol{(8x-5)^2}}$

\item L'épaisseur de $10$ feuilles est $60 \times 10^{-3}$ cm.\\
Pour $1\,000$ feuilles, il faut multiplier par $100$.\\
L'épaisseur de la pile en cm est donc :\\
$\begin{aligned}60 \times 100 \times 10^{-3}&=6\,000 \times 10^{-3}\\&=6 \times 10^{0} \\&=6\end{aligned}$\\
L'épaisseur de la pile est donc de ${\color{correction}\boldsymbol{6 \text{ cm}}}$.

\item On cherche le plus petit entier naturel $n$ vérifiant $h(n)>0$.\\
Comme $5n-38>0$ équivaut à $n>\dfrac{38}{5}$, le plus petit entier naturel $n$ est donc ${\color{correction}\boldsymbol{8}}$ car $\dfrac{38}{5} = 7,6$.

\item On met l'expression au même dénominateur :\\
$\begin{aligned}
\dfrac{1}{3}-\dfrac{2x+3}{x}&=\dfrac{x-3\times \left(2x+3\right)}{3x}\\
&=\dfrac{x -6x -9}{3x}\\
&=\dfrac{-5x -9}{3x}\\
\end{aligned}$\\
$\phantom{\dfrac{1}{3}-\dfrac{2x+3}{x}}=-\dfrac{5x +9}{3x}$

\end{enumerate}

\end{seanceD}

\end{document}